\subsection{4.4  Phase 4: Remediation Window and Check-In Cadence}

The remediation window is a bounded time period during which the APIP interventions are implemented and their effects are evaluated. A defined window serves two functions: it creates urgency for the remediation process, and it provides a structured decision point at which the APIP outcome is formally evaluated. The duration of the window should be calibrated to the intervention tier---Tier 1 interventions may warrant a 7--14 day window, while Tier 3 interventions may require 30--60 days to implement, validate, and stabilize. This mirrors the organizational PIP practice, where durations typically range from 30 to 90 days, calibrated to the complexity of the performance gap [8].

Check-ins are scheduled evaluation points within the window at which the agent's performance on APIP-relevant metrics is reviewed against the recovery trajectory. Following Kirkpatrick's [7] hierarchical evaluation principle, each check-in should assess improvement at multiple levels: immediate output quality (Level 1/2), behavioral pattern change on the failure-mode class (Level 3), and business outcome metrics (Level 4). An intervention that improves surface metrics without producing Level 3/4 improvement should be flagged as potentially superficial. Check-ins are early warning mechanisms: no improvement at the first check-in triggers tier escalation rather than waiting for window expiry. A recommended cadence is three check-ins at 25\%, 50\%, and 75\% of elapsed time.

